\section{Proofs for the intrinsic boundary atlas}
\label{sec:boundary-atlas-proofs}

We prove the assertions of
Section~\ref{sec:boundary-atlas-statement}.  All calculations are made
in the regular coefficient ring of a Schur chart.  Since the
identities are invariant under changes of frames, they descend by the
same argument as the Fitting identities in
Proposition~\ref{prop:full-fitting}.

\subsection{Universal determinant completion}

\begin{proof}[Proof of Theorem~\ref{thm:universal-det-completion}]
Put \(N=\dim\Sym^rV\) and \(s=\dim S\).  Choose a vector
space \(K\) of dimension \(N-s\) and a linear map
\(\rho:\Sym^rV\to K\) such that
\[
 G=(\gamma,\rho):\Sym^rV\xrightarrow{\sim}S\oplus K
\]
is an isomorphism.  For the generic endomorphism \(M\),
\[
 \det(\Sym^rM)=(\det M)^{\kappa(e,r)},\qquad
 \kappa(e,r)=\frac r e\binom{e+r-1}{r}
              =\binom{e+r-1}{r-1}.
\]
The eigenvalue identity proves the formula on diagonal matrices and
hence as a polynomial identity on \(\End(V)\).

Expand
\[
 \det(G\Sym^rM)=\det(G)d^{\kappa(e,r)}
\]
by Laplace expansion along the first \(s\) rows.  Every summand is
an \(s\)-minor of \(\gamma\Sym^rM\) times a complementary
minor of \(\rho\Sym^rM\).  Therefore the determinant belongs to
\(J_{\gamma,r}\).  Since \(\det G\ne0\),
\(d^{\kappa(e,r)}\in J_{\gamma,r}\).

If \(\gamma\) is an isomorphism, the ideal of maximal minors is
principal and generated by
\(\det(\gamma)\det(\Sym^rM)\), so the exponent is sharp.
\end{proof}

\subsection{The depth bound and graded multiplication}

\begin{proof}[Proof of Theorem~\ref{thm:finite-all-corank-depth}]
Apply Theorem~\ref{thm:universal-det-completion} with
\(e=4\), \(r=2\), \(\gamma=\gamma_R\), and
\(M=I_H\oplus T\).  Since \(\kappa(4,2)=5\),
\(d^5\in I_6(\gamma_R\Sym^2M)=J\).
Now \(\mathcal I_{\widehat D_R}=dJ\).  Therefore
\(d^6\in dJ\), so the nilradical generated by the class of \(d\)
satisfies \(N^6=0\).  Formula
\eqref{eq:global-graded-piece} then gives \(W_j=\varnothing\) for
\(j\ge6\).

For the multiplication statement work locally in
\(A=P/(dJ)\).  The class of \(rd^i\) in
\(N^i/N^{i+1}\) depends on \(r\) modulo
\[
 (d)+(J:d^{i-1}),
\]
and multiplication sends
\[
 (rd^i)(sd^j)=rs\,d^{i+j}.
\]
Because
\[
 (J:d^{i-1})+(J:d^{j-1})
 \subset (J:d^{i+j-1}),
\]
this factors through the canonical restriction to \(W_{i+j}\).
The image of \(1\otimes1\) is the generator \(d^{i+j}\), hence the
map is surjective.  The transition functions of \(d\) supply the
factor \(\mathcal L^{\otimes(i+j)}\), proving
\eqref{eq:atlas-multiplication}.
\end{proof}

\subsection{The first layer at projection corank two}

Assume throughout this subsection that \(A_H\) is injective.  After
eliminating its image, the residual ideal is the ideal of \(3\)-minors
of
\begin{equation}\label{eq:psi-beta-chi}
 \Psi_H(T)=
 \left[
   \beta_H(1_H\otimes T),\,
   \chi_H\Sym^2T
 \right],
 \qquad \dim E_H=3.
\end{equation}

\begin{lemma}[The contact symbol without mixed regularity]
\label{lem:contact-symbol-all-Ainj}
Modulo \((\delta)\), the residual ideal of
\eqref{eq:psi-beta-chi} is generated in degree four and equals
\[
 h_H(u)\Sym^4(\C^2_v)
\]
under \(T=uv^t\).  In particular
\[
 J+(\delta)=(\delta,g_0,\ldots,g_4)
\]
with the notation of
Theorem~\ref{thm:corank-two-mixed-kernel-atlas}.
\end{lemma}

\begin{proof}
On the rank-one cone write \(T=uv^t\).  The four mixed columns are
scalar multiples, with coefficients \(v_1,v_2\), of the two vectors
\[
 \beta_H(h_1\otimes u),\qquad
 \beta_H(h_2\otimes u),
\]
and the three quadratic columns are scalar multiples, with
coefficients \(v_1^2,v_1v_2,v_2^2\), of the single vector
\(\chi_H(u^2)\).  A \(3\)-minor containing three mixed columns
vanishes because the mixed span has dimension at most two.  A minor
containing at least two quadratic columns vanishes because those
columns are proportional.  Thus only minors with two mixed columns
and one quadratic column survive.  Their common \(u\)-factor is
\[
 \det[
   \beta_H(h_1\otimes u),
   \beta_H(h_2\otimes u),
   \chi_H(u^2)
 ],
\]
a binary quartic.  The five possible \(v\)-monomials span
\(\Sym^4(\C^2_v)\).  By the same determinant comparison used in
Lemma~\ref{lem:rootfree}, this binary quartic is the restriction
\(h_H\) of the Jacobian quartic to the pencil of hyperplanes
containing \(H\).  This proves the assertion without a rank
assumption on \(\beta_H\).
\end{proof}

The orbit classification in
Definition~\ref{def:mixed-kernel-orbit-type} is elementary projective
geometry.  We record normal forms because they are also used in the
colon calculation.  Order the basis of \(H\otimes W\) as
\[
 h_1w_1,\ h_1w_2,\ h_2w_1,\ h_2w_2.
\]
Up to \(GL(H)\times GL(W)\times GL(E_H)\), the mixed block can be
taken to be one of the following matrices:
\begin{equation}\label{eq:beta-normal-forms}
\begin{array}{c|c}
\text{type}&\beta\\ \hline
3,\ \text{off }Q&
\begin{psmallmatrix}0&1&0&0\\0&0&1&0\\1&0&0&-1\end{psmallmatrix}\\[3pt]
3,\ \text{on }Q&
\begin{psmallmatrix}0&1&0&0\\0&0&1&0\\0&0&0&1\end{psmallmatrix}\\[3pt]
2,\ \text{secant}&
\begin{psmallmatrix}0&1&0&0\\0&0&1&0\\0&0&0&0\end{psmallmatrix}\\[3pt]
2,\ \text{tangent}&
\begin{psmallmatrix}0&1&-1&0\\0&0&0&1\\0&0&0&0\end{psmallmatrix}\\[3pt]
2,\ H\text{-ruling}&
\begin{psmallmatrix}0&0&1&0\\0&0&0&1\\0&0&0&0\end{psmallmatrix}\\[3pt]
2,\ W\text{-ruling}&
\begin{psmallmatrix}0&1&0&0\\0&0&0&1\\0&0&0&0\end{psmallmatrix}\\[3pt]
1,\ \text{dual rank }2&
\begin{psmallmatrix}1&0&0&1\\0&0&0&0\\0&0&0&0\end{psmallmatrix}\\[3pt]
1,\ \text{dual rank }1&
\begin{psmallmatrix}1&0&0&0\\0&0&0&0\\0&0&0&0\end{psmallmatrix}\\[3pt]
0&0.
\end{array}
\end{equation}
For rank two, a projective kernel line meets the smooth quadric in
two reduced points, one double point, or is contained in one of the
two ruling families.  For rank one, the dual generator has tensor
rank two or one.  These alternatives are invariant and exhaustive,
proving the orbit statement.

\begin{lemma}[Exact transverse models for the nine mixed-kernel orbits]
\label{lem:nine-orbit-colons}
Let
\[
 P=\C[a,b,c,d],\qquad
 T=\begin{pmatrix}a&b\\c&d\end{pmatrix},\qquad
 \delta=ad-bc.
\]
For each normal form in \eqref{eq:beta-normal-forms} there is a
nonempty Zariski-open set of quadratic blocks \(\chi\) on which the
Hilbert functions of
\[
 P/\bigl((\delta)+(J:\delta^{j-1})\bigr)
\]
are those displayed in
Theorem~\ref{thm:corank-two-mixed-kernel-atlas}.
\end{lemma}

\begin{proof}
The two rank-three rows are
Theorem~\ref{thm:graded-nilpotent} and
Theorem~\ref{thm:decomposable-kernel-wall}.  For the remaining rows,
the assertion is a homogeneous determinant calculation.  We give
integer witnesses and the resulting colons; this simultaneously
proves nonemptiness of the relevant rank opens.

For the secant normal form take
\[
 \chi=
 \begin{pmatrix}
  1&2&-2\\-2&-1&1\\-2&1&-2
 \end{pmatrix}.
\]
A direct expansion of the \(3\)-minors followed by cancellation by
\(\delta\) gives
\[
\begin{split}
 (J:\delta)=(&c^3,c^2d,cd^2,d^3,\,
 a^2-c^2,ab-cd,b^2-d^2,\\
 &ac-2c^2,bc-2cd,ad-2cd,bd-2d^2),
\end{split}
\]
and \((J:\delta^2)=P\).  The first ideal contains \(\delta\), and
the quotient has Hilbert function \((1,4,3)\), hence length eight.

For the tangent normal form take
\[
 \chi=
 \begin{pmatrix}
  1&-2&-1\\0&1&0\\-1&-2&2
 \end{pmatrix}.
\]
Then
\[
\begin{split}
 (J:\delta)=(&c^3,c^2d,cd^2,d^3,\,
 a^2-10c^2,ab-10cd,b^2-10d^2,\\
 &ac+2c^2,bc+2cd,ad+2cd,bd+2d^2),
\end{split}
\]
again with \((J:\delta^2)=P\).  The same Hilbert function
\((1,4,3)\) follows.

For the \(H\)-ruling normal form take
\[
 \chi=
 \begin{pmatrix}
  1&-1&-1\\1&1&-1\\2&-1&0
 \end{pmatrix}.
\]
One obtains
\[
\begin{split}
 (J:\delta)=(&c^4,c^3d,c^2d^2,cd^3,d^4,\,
 5ac^2-c^3,5bc^2-c^2d,\\
 &5acd-c^2d,5bcd-cd^2,
 5ad^2-cd^2,5bd^2-d^3,\\
 &a^2-ac,\,2ab-ad-bc,\,b^2-bd),
\end{split}
\]
while
\[
 (J:\delta^2)=(a,b,c,d).
\]
Here \(\delta\notin(J:\delta)\), so \(\delta^2\notin J\);
on the other hand \(\delta\in(J:\delta^2)\), so
\(\delta^3\in J\).  After adjoining \(\delta\), the second-layer
quotient has Hilbert function
\[
 (1,4,6,4)
\]
and length fifteen, while the third layer has length one.

For the \(W\)-ruling normal form take
\[
 \chi=
 \begin{pmatrix}
  1&-1&1\\2&1&2\\1&0&-2
 \end{pmatrix}.
\]
Then
\[
 (J:\delta)=
 (ac,bc,ad,bd,c^3,c^2d,cd^2,d^3),
 \qquad
 (J:\delta^2)=P.
\]
This ideal already contains \(\delta\).  The Hilbert function of the
second layer is
\[
 1,4,6,4,5,6,7,8,\ldots,
\]
hence its Hilbert polynomial is \(n+1\) from degree four onward.

For the dual-rank-two rank-one normal form take
\[
 \chi=
 \begin{pmatrix}
 -2&1&-1\\0&0&-2\\-2&0&-1
 \end{pmatrix}.
\]
After adjoining \(\delta\), a homogeneous basis for the second colon
may be chosen as
\[
 (\delta,\,
 a^2c,ac^2,a^2d,abd,b^2d,acd,ad^2,bd^2).
\]
Its Hilbert function is
\[
 1,4,9,8,10,12,14,16,\ldots,
\]
so the Hilbert polynomial is \(2n+2\) from degree three onward.
Moreover
\[
 (J:\delta^2)=(a,b,c,d),
\]
and the third layer is the reduced vertex.

For the dual-rank-one normal form take
\[
 \chi=
 \begin{pmatrix}
 -2&2&-1\\1&1&1\\2&1&-2
 \end{pmatrix}.
\]
A homogeneous basis after adjoining \(\delta\) is
\[
\begin{split}
(\delta,\,
&a^3+4a^2c+3ac^2,\,
 a^2b+4a^2d+3acd,\\
&ab^2+4abd+3ad^2,\,
 b^3+4b^2d+3bd^2),
\end{split}
\]
with Hilbert function
\[
 1,4,9,12,15,18,21,\ldots.
\]
The next colon is, after an invertible linear change,
\[
 (c^2,cd,d^2,a,b),
\]
so the third layer has Hilbert function \((1,2)\) and length three.

Finally, when \(\beta=0\), surjectivity of
\([\beta,\chi]\) forces \(\chi\) to be invertible.  Hence
\[
 J=I_3(\chi\Sym^2T)
   =(\det\chi)\det(\Sym^2T)
   =(\delta^3),
\]
because \(\det(\Sym^2T)=\delta^3\) in dimension two.  This proves the
last row exactly.

All displayed Hilbert functions are computed from the displayed
homogeneous bases by counting standard monomials.  The role of these
integral points is only to provide convenient base points.
Genericity is proved separately in
Theorem~\ref{thm:generic-slice-certification}: the stabilizer orbit
of each normal form is supplemented by an explicit transverse affine
slice, and the residual colons are computed over the rational
function field of that slice.  The resulting initial ideals give
exactly the Hilbert functions displayed above.  Hence the open
conditions in the statement are genuine dense-open conditions on the
whole orbit family, not constant-rank loci inferred from one fibre.
\end{proof}

\begin{proof}[Proof of Theorem~\ref{thm:corank-two-mixed-kernel-atlas}]
Lemma~\ref{lem:contact-symbol-all-Ainj} gives the first layer and its
contact interpretation.  The orbit list follows from
\eqref{eq:beta-normal-forms}.  Lemma~\ref{lem:nine-orbit-colons}
gives the second and third layers on each coefficient-transverse open.
The nilpotency index is
\[
 1+\min\{k:\delta^k\in J\}
\]
by Theorem~\ref{thm:global-residual-colon}; the displayed colons give
the values in the table.

If \(b\le1\), two mixed columns never span two independent target
directions, so the determinant in the proof of
Lemma~\ref{lem:contact-symbol-all-Ainj} vanishes identically and
\(h_H=0\).  If the kernel is the \(H\)-ruling
\(h_0\otimes W\), one of the two mixed vectors vanishes for every
\(u\), so again \(h_H=0\).  For the \(W\)-ruling, the mixed span
drops only at \(u=w_0\), giving the forced contact root but not
identical vanishing on the transverse open.  This proves the
remaining assertions.
\end{proof}

\subsection{Rank-drop depth and the finite coefficient atlas}

\begin{proof}[Proof of Proposition~\ref{prop:rankdrop-coranktwo-depth}]
Eliminate the rank-\(a\) block \(A_H\) and put \(q=6-a\).
The remaining target has dimension \(q\).  A maximal residual minor
uses at most \(b\) independent mixed columns, each of transverse
degree one.  Every remaining column must come from the quadratic
block and has transverse degree two.  Therefore every nonzero
generator of \(J\) has degree at least
\[
 b+2(q-b)=2q-b.
\]
If \(2k<2q-b\), the homogeneous polynomial \(\delta^k\) cannot lie in
\(J\), proving
\[
 k\ge q-\lfloor b/2\rfloor.
\]
Theorem~\ref{thm:finite-all-corank-depth} gives \(k\le5\).

The quadratic block has only three source dimensions when \(c=2\).
Surjectivity of the original quotient after eliminating \(A_H\)
therefore implies \(b+3\ge q\).  Combining this with \(b\le4\) gives
exactly the table in the statement.  The equivalence between
\(d^k\in J\) and the augmented-rank condition
\(\varepsilon_k(J)=0\) is ordinary linear algebra in the degree
\(2k\) homogeneous component.  In the type \((a,b)=(0,3)\) the lower
bound already gives \(k\ge5\), while \(d^5\in J\), so \(k=5\)
exactly.
\end{proof}

The same coefficient construction works for \(c=3,4\).  Because
\(d^5\in J\), it terminates at \(k=5\).  Colon ideals are kernels of
the maps
\[
 (P_c)_m\xrightarrow{\ \cdot d^{j-1}\ }
 (P_c)_{m+c(j-1)}/J_{m+c(j-1)},
\]
hence are obtained by finite syzygy elimination.  Noetherianity gives
finite homogeneous generating sets.  Rank stratification of the
resulting coefficient matrices, followed by the usual flattening
stratification of the finitely presented graded quotients, gives the
finite atlas in Definition~\ref{def:macaulay-atlas}.  This is the
finite algebraic replacement for an unspecified infinite colon tower.

\subsection{Projection corank three}

\begin{lemma}[Functorial symbol containment under mixed-rank drop]
\label{lem:corank3-functorial-symbol}
Fix the rank of \(A_H\) at projection corank three.  In every
homogeneous degree, the right \(GL(W)\)-module generated by the
residual maximal minors for a specialization of the mixed block is a
quotient of the corresponding universal maximal-mixed-rank symbol
module.  After multiplication by coordinate functions the same is
true degree by degree.
\end{lemma}

\begin{proof}
After eliminating \(\im A_H\), a residual minor with \(u\)
mixed columns and \(q-u\) quadratic columns is the image of the
natural alternating map
\[
 \bigwedge^u(H\otimes W)\otimes
 \bigwedge^{q-u}\Sym^2W
 \longrightarrow \bigwedge^q E_H.
\]
The construction is polynomial and functorial in the mixed map
\(H\otimes W\to E_H\).  Replacing that map by one of smaller
rank factors the displayed map through the corresponding quotient of
\(H\otimes W\); it cannot create a new right
\(GL(W)\)-constituent.  Tensoring with homogeneous coordinate
functions and taking images preserves this quotient relation.
\end{proof}
\begin{lemma}[No cubic determinant power at corank three]
\label{lem:no-d3-corank3}
At projection corank three, \(d^3\notin J\).
\end{lemma}

\begin{proof}
Here \(\dim W=3\) and \(\deg d=3\).  First suppose
\(A_H\) has rank one.  After eliminating its image the residual target
has dimension five.  The maximal-rank mixed case has
\(\rank\overline B_H=3\).  By
Lemma~\ref{lem:corank3-functorial-symbol}, every mixed-rank drop is a
quotient of the symbol module treated below, degree by degree.  Thus a
determinant representation absent from the universal symbol cannot be
created by specialization.

The degree-seven, degree-eight and degree-nine symbols of the
five-minors use respectively three, two and one mixed columns.  As
right \(GL(W)\)-representations they are contained in
\[
 \bigwedge^3W\otimes\bigwedge^2\Sym^2W,\qquad
 \bigwedge^2W\otimes\bigwedge^3\Sym^2W,\qquad
 W\otimes\bigwedge^4\Sym^2W.
\]
The elementary Schur decompositions are
\[
 \bigwedge^2\Sym^2W=\mathbb S_{(3,1,0)}W,
\]
\[
 \bigwedge^3\Sym^2W=
 \mathbb S_{(4,1,1)}W\oplus
 \mathbb S_{(3,3,0)}W,
\]
and
\[
 \bigwedge^4\Sym^2W=\mathbb S_{(4,3,1)}W.
\]
Consequently the degree-eight symbol is a sum of
\[
 \mathbb S_{(5,2,1)},\
 \mathbb S_{(4,4,0)},\
 \mathbb S_{(4,3,1)},\
 \mathbb S_{(4,2,2)},
\]
and the degree-nine symbol is a sum of
\[
 \mathbb S_{(5,3,1)},\
 \mathbb S_{(4,4,1)},\
 \mathbb S_{(4,3,2)}.
\]
The degree-seven symbol is
\(\mathbb S_{(4,2,1)}\).
To obtain a degree-nine element of the ideal one may multiply the
degree-seven symbol by degree-two coordinate functions, the
degree-eight symbol by degree-one coordinate functions, or take the
degree-nine symbol itself.  None can contain
\[
 (\det W)^3=\mathbb S_{(3,3,3)}W:
\]
every displayed highest weight has first part at least four, and the
Littlewood--Richardson rule only adds boxes.  Since \(d^3\) spans
this determinant representation, it is not in \(J\).

If \(A_H=0\), the residual target has dimension six.  A degree-nine
minor can occur only when the mixed block has rank three and uses all
three mixed columns and three quadratic columns.  Its right type is
contained in
\[
 \det W\otimes\bigwedge^3\Sym^2W
 =
 \mathbb S_{(5,2,2)}W
 \oplus
 \mathbb S_{(4,4,1)}W,
\]
again with no \((3,3,3)\) summand.  If the mixed rank drops, the
minimal degree is larger than nine.  This proves the assertion in all
cases.
\end{proof}

\begin{proof}[Proof of Theorem~\ref{thm:corank-three-depth}]
Lemma~\ref{lem:no-d3-corank3} gives \(d^3\notin J\), while
Theorem~\ref{thm:finite-all-corank-depth} gives \(d^5\in J\).
Therefore the first determinant power in \(J\) is \(d^4\) or \(d^5\).
Definition~\ref{def:macaulay-atlas} says precisely that
\(d^4\in J\) when \(\varepsilon_4(J)=0\).  The nilpotency index is
one plus this first exponent.
\end{proof}

\subsection{Projection corank four}

We use one standard plethysm calculation (see, for example,
\cite{FultonHarris} for Schur functors and character methods).  The
identity may also be checked directly by the Weyl character formula:
\begin{equation}\label{eq:plethysm-corank4}
 \bigwedge\nolimits^6\Sym^2V
 =
 \mathbb S_{(5,4,2,1)}V
 \oplus
 \mathbb S_{(4,4,4,0)}V.
\end{equation}
Both summands occur with multiplicity one.

\begin{lemma}[Effective Cauchy--Pieri multiplication in the matrix ring]
\label{lem:effective-pieri-matrix-ring}
Let \(A=\Sym(V^*\otimes V)=\C[\End(V)]\).  In the Cauchy
decomposition of \(A\), put
\[
 \lambda=(4,4,4,0),\qquad \mu=(4,0,0,0),\qquad
 \nu=(4,4,4,4).
\]
The actual polynomial multiplication
\[
 A_{12}[\lambda]\otimes A_4[\mu]\longrightarrow A_{16}
\]
has nonzero projection to \(A_{16}[\nu]\).  This projection has
multiplicity one on both the left and right Schur factors and, after
the identifications
\[
 \mathbb S_\lambda V\simeq
 (\det V)^4\otimes(\Sym^4V)^*,\qquad
 \mathbb S_\mu V=\Sym^4V,
\]
is a nonzero scalar multiple of the evaluation pairing.
\end{lemma}

\begin{proof}
The skew diagram \(\nu/\lambda\) consists of four boxes, one in
each column, hence is a horizontal four-strip.  Pieri's rule gives
\(c^{\nu}_{\lambda,\mu}=1\).  What is needed here is
nonvanishing for the multiplication of polynomial functions, not
merely the existence of an abstract equivariant map.

Use the standard bideterminant realization of the Cauchy summands in
\(\Sym(V^*\otimes V)\).  Multiply a standard
\(\lambda\)-bideterminant by the row-shaped
\(\mu\)-bideterminant which fills the fourth row.  Straightening
contains the standard \(\nu\)-bideterminant with coefficient
\(1\).  Therefore the \(\nu\)-projection of the actual
coordinate-ring product is nonzero.  Since the relevant
Littlewood--Richardson multiplicity is one, Schur's lemma identifies
the two factor maps with the unique Pieri maps up to nonzero scalar.
For the displayed determinant twist those maps are the perfect
evaluation pairings.
\end{proof}
\begin{lemma}[The Jacobian projection]
\label{lem:jacobian-projection}
Under
\[
 \bigwedge\nolimits^6\Sym^2V
 \simeq
 (\det V)^5\otimes
 \left(\bigwedge\nolimits^4\Sym^2V\right)^*,
\]
the summand
\(\mathbb S_{(4,4,4,0)}V\) is dual, up to the displayed determinant
twist, to
\[
 \det V\otimes\Sym^4V\subset
 \bigwedge\nolimits^4\Sym^2V.
\]
The projection
\[
 \bigwedge\nolimits^4\Sym^2V
 \longrightarrow
 \det V\otimes\Sym^4V
\]
sends \(q_1\wedge\cdots\wedge q_4\), up to a nonzero scalar, to
\[
 \det\left(\frac{\partial q_i}{\partial x_j}\right).
\]
\end{lemma}

\begin{proof}
Duality of Schur functors gives
\[
 (\det V)^5\otimes
 \mathbb S_{(4,4,4,0)}V^*
 \simeq
 \det V\otimes\Sym^4V.
\]
The Jacobian determinant is alternating in the four quadrics and is
\(GL(V)\)-equivariant with exactly this target weight.  It is nonzero,
for instance on
\(x_1^2,x_2^2,x_3^2,x_4^2\).  Multiplicity one in
\eqref{eq:plethysm-corank4} therefore identifies it with the stated
projection.
\end{proof}

\begin{proof}[Proof of Theorem~\ref{thm:corank-four-jacobian}]
At \(H=0\) the ideal \(J\) is generated in degree twelve by the
\(6\)-minors of \(\gamma_R\Sym^2T\).  Under right multiplication of
\(T\), their span is a quotient of the representation
\(\bigwedge^6\Sym^2V\).  The polynomial \(d^3\) spans
\((\det V)^3=\mathbb S_{(3,3,3,3)}V\), which does not occur in
\eqref{eq:plethysm-corank4}.  Hence \(d^3\notin J\).

Let
\[
 \lambda_R=\bigwedge\nolimits^6\gamma_R.
\]
By Lemma~\ref{lem:jacobian-projection}, the component of
\(\lambda_R\) on the
\(\mathbb S_{(4,4,4,0)}\)-summand is nonzero exactly when the
Jacobian quartic \(F_R\) is nonzero.  For \(R\in G_4^\circ\), the
quartic \(Y_R=V(F_R)\) is smooth, so this component is nonzero.

Let \(A=\C[\End(V)]\).  The nonzero Jacobian component just
identified gives a nonzero vector in the
\(\lambda=(4,4,4,0)\) Cauchy component of \(J_{12}\).
By Lemma~\ref{lem:effective-pieri-matrix-ring}, multiplication by
the \(\mu=(4)\) component of \(A_4\) has a nonzero
\(\nu=(4,4,4,4)\)-projection.  Because the induced pairing
\[
 \mathbb S_{(4,4,4,0)}V\otimes\Sym^4V
 \longrightarrow(\det V)^4
\]
is nondegenerate, one may choose the degree-four matrix coefficient
to pair nontrivially with the specific nonzero Jacobian vector.

The \(\nu\)-Cauchy component of \(A_{16}\) is one-dimensional
on both sides and is the line generated by the fourth power of the
matrix determinant.  The product belongs to the ideal \(J\), so
its nonzero \(\nu\)-projection shows that this line meets \(J\)
nontrivially.  Hence the generator itself belongs to the ideal:
\[
 d^4\in J.
\]
This uses the nonzero structure constant of actual polynomial
multiplication supplied by bideterminant straightening; representation
occurrence alone is not being used as a membership argument.  Together with \(d^3\notin J\), this
gives nilpotency index five.
\end{proof}

\subsection{Specialization and associated supports}

The coefficient matrices above have polynomial entries, but this fact
alone does not justify fibrewise statements about embedded associated
points or base change of colons.  The required order of argument is
carried out in Section~\ref{sec:relative-primary-specialization}.
There we first flatten the finite presentations controlling each
kernel \((J:d^{j-1})\) and its cokernel, then apply relative
assassins, generic freeness at the associated generic points, and
Noetherian induction.  Proposition
\ref{prop:associated-prime-refinement} follows from that argument.

The same section computes the generic corank-two primary signatures
over rational function fields and gives explicit one-parameter orbit
degenerations.  Thus the specialization statements used by the paper
are consequences of exact family ideals and relative associated-point
theory, not of generic flatness alone.
